% ============================================================
%  Section 2 -- Description conceptuelle du probleme
% ============================================================

\chapter{Description conceptuelle du problème}

\section{Nature d'un ticket de support IT}

Un ticket de support IT est une demande courte, souvent ambigue, qui melange un
symptome utilisateur, un contexte applicatif et parfois une urgence implicite.
Dans le jeu de donnees du projet, les exemples couvrent des problemes de VPN,
Outlook, BitLocker, mot de passe, lenteur poste de travail, imprimante ou acces
reseau. Le fichier \file{data/inputs/benchmark/tickets_scenario_1.csv} contient
40 tickets exploites par les scripts de benchmark, dont 10 sont utilises dans la
campagne principale enregistree.

Une reponse utile doit satisfaire plusieurs criteres simultanement :

\begin{itemize}
  \item identifier le probleme probable sans surinterpreter le ticket ;
  \item proposer des actions concretes et ordonnees ;
  \item exploiter les sources internes disponibles lorsque la question depend
  d'une procedure ou d'un statut ;
  \item eviter les hallucinations, en particulier les numeros, chemins, etats de
  service ou politiques internes inventes ;
  \item respecter les contraintes de securite et de confidentialite.
\end{itemize}

Ces criteres sont difficiles a evaluer manuellement a grande echelle. C'est le
role de \bibopsname{} : transformer une comparaison informelle entre agents en
campagne instrumentee, rejouable et mesurable.

\section{Limites d'un LLM unique}

L'architecture \emph{LLM Unique} correspond a un modele appele directement sur
le ticket, sans acces outil. Elle est simple a deployer, mais presente trois
limites structurelles.

\subsection{Absence d'etat externe}

Un LLM sans outil ne peut pas savoir si un service est indisponible, si une
procedure interne a change ou si un incident en cours existe. Il peut memoriser
des connaissances generales, mais il ne peut pas verifier un etat local comme la
table SQLite \code{serveurs_it}. Pour un ticket du type « Est-ce que le service
Mail est en ligne ? », cette absence d'acces a l'etat rend la reponse
necessairement speculative.

\subsection{Risque d'hallucination procedural}

Les tickets de support invitent le modele a produire des listes d'etapes. Or une
liste detaillee peut etre fausse tout en semblant professionnelle. Le risque est
eleve lorsque la reponse inclut des URLs internes, numeros de support, noms de
serveurs ou commandes d'administration. Le benchmark fact-checking du depot
illustre ce probleme : l'affirmation « le support IT Michelin est joignable au
1234 » est traitee comme non verifiable par l'agent fact-checker, alors que des
reponses generatives peuvent reprendre ce type de valeur comme si elle etait
certaine.

\subsection{Evaluation trompeuse par fluidite}

Un modele unique peut produire une reponse longue et bien structuree, mais la
qualite apparente ne suffit pas. L'evaluation doit separer la lisibilite, la
pertinence, la fidelite aux sources et l'absence de risque. Cette separation est
assuree dans le projet par un juge LLM pour la qualite, un inspecteur de
securite deterministe et des metriques operationnelles.

\section{Interet et complexite d'un systeme agentique}

L'architecture multi-agents introduit des outils et une boucle de decision. Le
LLM ne doit plus seulement repondre : il doit choisir s'il a besoin d'une source,
appeler l'outil pertinent, integrer le resultat, puis formuler une reponse
finale. Cette approche est proche du paradigme ReAct, qui couple raisonnement et
action \cite{react2023}.

Elle resout certains problemes du LLM unique, mais en cree d'autres.

\begin{itemize}
  \item \textbf{Routage.} Le modele peut choisir le mauvais outil ou appeler
  plusieurs fois le meme outil.
  \item \textbf{Robustesse.} Un outil peut etre lent, indisponible ou retourner
  un resultat vide.
  \item \textbf{Securite.} Les donnees recuperees par RAG ou API peuvent
  contenir des instructions malicieuses.
  \item \textbf{Traçabilite.} L'agent doit conserver les appels outils pour
  expliquer sa reponse et auditer les echecs.
\end{itemize}

Le depot traite ces risques par des politiques d'outils, des traces JSONL, une
normalisation des arguments et un blocage des appels dupliques. Ces protections
ne rendent pas l'agent parfait, mais elles transforment les erreurs en signaux
observables.

\section{Probleme d'evaluation}

Comparer deux architectures LLM est difficile car il n'existe pas une seule
metrique universelle. Un score de qualite eleve peut masquer un cout trop grand,
une latence excessive ou une faille de securite. Inversement, une architecture
rapide et peu couteuse peut etre inutilisable si ses reponses ne passent pas un
seuil minimal.

\bibopsname{} formalise donc l'evaluation comme un probleme multicritere. La
politique composite du depot utilise les poids suivants :

\begin{table}[H]
\centering
\begin{tabularx}{0.78\textwidth}{Yc}
\toprule
\textbf{Dimension} & \textbf{Poids} \\
\midrule
Qualite & 40 \% \\
Securite & 35 \% \\
FinOps / cout & 10 \% \\
Latence & 10 \% \\
GreenOps & 5 \% \\
\bottomrule
\end{tabularx}
\caption{Ponderation de \code{CompositePolicy}.}
\label{tab:composite-weights}
\end{table}

Le score composite n'est pas seul decisionnaire. Des gates bloquants imposent
notamment un score qualite minimal de 7/10 et un score securite minimal de 6/10.
Cette conception est volontairement stricte : elle evite qu'une architecture
faible en qualite soit declaree deployable uniquement parce qu'elle est rapide
ou bon marche.

\section{Risques specifiques aux agents connectes}

L'ajout d'outils et de flux externes augmente la surface d'attaque. Un agent qui
lit une documentation, un flux SSE ou une API peut recevoir du contenu qui
ressemble a des instructions systeme. Les attaques de prompt injection exploitent
precisement cette ambiguite entre donnees et instructions \cite{promptinjection2023}.

La Racing Arena transpose ce risque dans une simulation F1. Les equipes doivent
prendre des decisions de pit-stop a partir d'une telemetrie diffusee par un hub.
Une equipe attaquante, Team Psi, tente d'influencer les autres via des messages
directs, des usurpations d'autorite ou du RAG empoisonne. Le probleme conceptuel
devient alors : comment evaluer la robustesse d'un agent non seulement sur une
reponse statique, mais dans une boucle de decision continue, exposee a des
messages concurrents et a des APIs faibles ?

\section{Contraintes de conception}

Le projet impose plusieurs contraintes d'ingenierie.

\begin{itemize}
  \item \textbf{Executabilite locale.} Les tests unitaires doivent passer sans
  appeler de vrais LLM ; les appels reseau sont mocks ou isoles.
  \item \textbf{Observabilite.} Chaque benchmark doit produire des artefacts
  persistants : JSON, graphiques, traces ou rapports.
  \item \textbf{Extensibilite.} Les evaluateurs et adaptateurs doivent pouvoir
  etre remplaces ou ajoutes sans modifier tout le pipeline.
  \item \textbf{Reproductibilite.} Les scripts CLI exposent des commandes
  documentees pour relancer les experiences.
  \item \textbf{Pragmatisme.} Les outils restent simples : SQLite pour l'etat de
  service, JSON pour la base de connaissances, ChromaDB pour le RAG.
\end{itemize}

Ces contraintes expliquent le choix d'un framework Python modulaire, testable,
avec une CLI Typer et des scripts de developpement qui peuvent etre lances
independamment.
